% Irodov Part III §3.4 problems 176--188, source PDF pp.125--128,
% printed pp.116--119; figures 3.40--3.51 visually inspected.
\electrocurrentproblem{176}{116}
图3.40：$N$ 只电源同向串联成闭合回路，
每只电源的电动势 $\mathcal E_i$
正比于其内阻 $R_i$，
$\mathcal E_i=\alpha R_i$，
导线电阻忽略。
求 (a) 环路电流；(b) 任意两节点 A、B
之间的电势差，它们把回路分成 $n$ 与
$N-n$ 段。
\exsolution
依 Kirchhoff 第二定律，
全回路总电动势
$\sum_{i=1}^{N}\mathcal E_i
=\alpha\sum_{i=1}^{N}R_i$，
总电阻为 $\sum_iR_i$，
所以
\[
\boxed{I=
\frac{\sum_i\mathcal E_i}{\sum_iR_i}
=\alpha}.
\]
沿任一串联源小段跨过其正向电动势时，
电势增加 $\mathcal E_i$，
流过内阻下降 $IR_i$。
因 $\mathcal E_i=IR_i$，
每一段的净电势跃变都为零，
故任意两连接节点电势相等：
\[
\boxed{\varphi_A-\varphi_B=0}.
\]
\excheck 非零电流 $I=\alpha$
与各段端电压为零并不矛盾：
电源电动势恰好在每段内阻上消耗。

\electrocurrentproblem{177}{116--117}
图3.41，两个理想电源
$\mathcal E_1=1.0\,\mathrm V$、
$\mathcal E_2=2.5\,\mathrm V$，
电阻
$R_1=10\,\Omega$、
$R_2=20\,\Omega$，
电容器 C 两极板接于点 A、B。
求平衡时 $\varphi_A-\varphi_B$。
\exsolution
稳定直流时电容器断路，
并不改变外部闭合回路的电流。
按图中两电源反向串联极性，
电流
\[
I=\frac{\mathcal E_1-\mathcal E_2}{R_1+R_2}
=\frac{1.0-2.5}{10+20}
=-0.050\,\mathrm A
\]
（负号表明实际电流与所取正向相反）。
由 A、B 两节点间的 $R_1$ 支路电势差，
\[
\varphi_A-\varphi_B=IR_1
=\frac{(\mathcal E_1-\mathcal E_2)R_1}
{R_1+R_2}.
\]
所以
\[
\boxed{\varphi_A-\varphi_B=-0.50\,\mathrm V}.
\]
\excheck 最终 $C$ 不导通直流，
但其两板仍可维持非零电势差。

\electrocurrentproblem{178}{117}
图3.42：电源电动势
$\mathcal E=5.0\,\mathrm V$，
内阻 $R=0.10\,\Omega$，
外接并联电阻
$R_1=4.0\,\Omega$、
$R_2=6.0\,\Omega$。
求各电阻上的电流。
\exsolution
外电路等效电阻
$R_L=R_1R_2/(R_1+R_2)$。
总电流
$I=\mathcal E/(R+R_L)$，
外电路端电压
$U=IR_L$。
由并联分流：
\[
I_1=\frac U{R_1}
=\boxed{\frac{\mathcal E R_2}
{R(R_1+R_2)+R_1R_2}},
\]
\[
I_2=\frac U{R_2}
=\boxed{\frac{\mathcal E R_1}
{R(R_1+R_2)+R_1R_2}}.
\]
代入数值：
分母 $0.10(10)+24=25$，
得到
\[
\boxed{I_1=1.2\,\mathrm A,\qquad
I_2=0.80\,\mathrm A}.
\]
\excheck 总电流为2 A，
内阻压降0.20 V，
外电路端电压4.8 V；
$1.2\times4=0.8\times6=4.8$，相符。

\electrocurrentproblem{179}{117}
图3.43：长为 $\ell$、总电阻
$R_0$ 的均匀电位器两端维持电压
$V_0$，滑动端距参考端为 $x$。
负载电阻 $R$ 接在滑片和该参考端之间。
求负载电压 $V(x)$，讨论 $R\gg R_0$。
\exsolution
电位器靠参考端一段电阻
$R_{\rm low}=R_0x/\ell$，
其余一段
$R_{\rm high}=R_0(\ell-x)/\ell$。
负载 $R$ 与 $R_{\rm low}$ 并联，
二者组合与 $R_{\rm high}$ 串联。
分压：
\[
V=V_0\frac{R\parallel R_{\rm low}}
{R_{\rm high}+R\parallel R_{\rm low}}
=V_0\frac{RR_{\rm low}}
{R(R_{\rm high}+R_{\rm low})
+R_{\rm high}R_{\rm low}}.
\]
代入分段电阻：
\[
\boxed{V(x)=
V_0\frac{Rx}
{R\ell+R_0x(\ell-x)/\ell}}.
\]
当 $R\gg R_0$，负载影响可忽略：
\[
\boxed{V(x)\simeq V_0\,\frac{x}{\ell}}.
\]
\excheck $x=0$ 时 $V=0$，
$x=\ell$ 时 $V=V_0$；
中间段的真实电压一般小于理想线性分压。

\electrocurrentproblem{180}{117}
两只电源各具有电动势
$\mathcal E_1,\mathcal E_2$
及内阻 $R_1,R_2$，
以相同端子极性并联。
求该并联源的等效电动势
及内阻（戴维南等效）。
\exsolution
设两电源并联输出端电压为 $U$，
各支路向外输出的电流
$I_i=(\mathcal E_i-U)/R_i$，
合电流
\[
I=\frac{\mathcal E_1}{R_1}
+\frac{\mathcal E_2}{R_2}
-U\left(\frac1{R_1}+\frac1{R_2}\right).
\]
把它改写为单个戴维南源形式
$I=(\mathcal E_{\rm eq}-U)/R_{\rm eq}$，
通过比较斜率与截距：
\[
\boxed{R_{\rm eq}=\frac{R_1R_2}{R_1+R_2}},
\]
\[
\boxed{\mathcal E_{\rm eq}=
\frac{\mathcal E_1R_2+\mathcal E_2R_1}
{R_1+R_2}}.
\]
\excheck 两电源即使未接外负载，
在 $\mathcal E_1\ne\mathcal E_2$ 时
内部可存在环流；上式给的是外端口等效。

\electrocurrentproblem{181}{117}
图3.44：两并行电源支路，
$\mathcal E_1=1.5\,\mathrm V$、
$\mathcal E_2=3.7\,\mathrm V$，
对应电阻
$R_1=10\,\Omega$、
$R_2=20\,\Omega$，
另有连接电阻 $R=5.0\,\Omega$。
求通过 $R$ 的电流大小与方向。
\exsolution
将上、中的两源电阻支路
对底部电阻 $R$ 看作戴维南等效。
依图中两电动势极性，
两端开路电势差
\[
\mathcal E_{\rm th}=
\frac{R_1\mathcal E_2-R_2\mathcal E_1}
{R_1+R_2},
\]
内部等效电阻
$R_{\rm th}=R_1R_2/(R_1+R_2)$。
故通过底部连接电阻的电流：
\[
\boxed{I=
\frac{R_1\mathcal E_2-R_2\mathcal E_1}
{R(R_1+R_2)+R_1R_2}}.
\]
代入：
\[
I=\frac{10(3.7)-20(1.5)}
{5(10+20)+10(20)}
=\frac7{350}
=\boxed{0.020\,\mathrm A}.
\]
正号对应原书图3.44
\emph{自左向右}流过底部 $R$。
\excheck 如果只看电动势
$\mathcal E_2>\mathcal E_1$
而忽略内串电阻，
不能正确判定流向。

\electrocurrentproblem{182}{117}
图3.45：A、B 两节点间有三条并联支路，
分别有 $(\mathcal E_1,R_1)$、
$(\mathcal E_2,R_2)$、
$(\mathcal E_3,R_3)$，
其中第三电源按图反向连接。
数据：
$\mathcal E_1=1.5\,\mathrm V$、
$\mathcal E_2=2.0\,\mathrm V$、
$\mathcal E_3=2.5\,\mathrm V$，
$R_1=10\,\Omega$、
$R_2=20\,\Omega$、
$R_3=30\,\Omega$。
求 (a) $R_1$ 电流；(b) $\varphi_A-\varphi_B$。
\exsolution
设 $U=\varphi_A-\varphi_B$，
三支路自 A 至 B 的电流分别取
\[
I_1=\frac{U-\mathcal E_1}{R_1},\quad
I_2=\frac{U-\mathcal E_2}{R_2},\quad
I_3=\frac{U+\mathcal E_3}{R_3}.
\]
A、B 之间无外部注入电流，
所以 $I_1+I_2+I_3=0$。
故
\[
\boxed{U=
\frac{\mathcal E_1/R_1+
\mathcal E_2/R_2-\mathcal E_3/R_3}
{1/R_1+1/R_2+1/R_3}}.
\]
数据代入：
$U=(0.15+0.10-0.08333)/
(0.1+0.05+0.03333)\,\mathrm V
\simeq0.909\,\mathrm V$，
因此
\[
\boxed{\varphi_A-\varphi_B\simeq0.90\,\mathrm V}.
\]
经过 $R_1$ 向上（B 到 A）的电流
\[
I_{1,\uparrow}=
\frac{\mathcal E_1-U}{R_1}
\simeq\boxed{0.059\,\mathrm A\approx0.060\,\mathrm A}.
\]
整理一般式：
\[
\boxed{I_{1,\uparrow}=
\frac{R_3(\mathcal E_1-\mathcal E_2)
+R_2(\mathcal E_1+\mathcal E_3)}
{R_1R_2+R_2R_3+R_3R_1}}.
\]
\excheck 书末答案的纯文本提取似把
第二项 $\mathcal E_3$ 混成 $\mathcal E_2$；
用题给数据和书末 $0.06$ A
可独立核实应为 $\mathcal E_3$。

\electrocurrentproblem{183}{118}
图3.46：电源 $\mathcal E_0$、$\mathcal E$
和电阻 $R_1,R_2,R_3,R$
按图连接，电源内阻可忽略。
求通过右下方电阻 $R$ 的电流。
\exsolution
图中左侧理想电源 $\mathcal E_0$
把其上下节点电势差固定，
与它并联的 $R_1$ 仅改变电源输出电流，
不改变流过 $R$ 的求解。
把左侧 $\mathcal E_0$ 与 $R_2,R_3$
看作对右侧支路的戴维南源：
开路电压由 $R_2,R_3$ 的分压决定，
\[
U_{\rm th}=
\frac{R_3}{R_2+R_3}\mathcal E_0.
\]
关闭理想源时，内部等效电阻为
$R_{\rm th}=R_2R_3/(R_2+R_3)$。
右侧电源 $\mathcal E$ 同向作用，
故
\[
I=\frac{\mathcal E+U_{\rm th}}{R+R_{\rm th}}
=\boxed{
\frac{\mathcal E(R_2+R_3)+\mathcal E_0R_3}
{R(R_2+R_3)+R_2R_3}}.
\]
\excheck 正如原书简答，
结果\emph{不含 $R_1$}，
因为它并联在理想电源两端，
不会改变左右外端口的戴维南电压。

\electrocurrentproblem{184}{118}
图3.47：由
$R_1=10\,\Omega$、
$R_2=20\,\Omega$、
$R_3=30\,\Omega$ 及电源
$\mathcal E_1=4.0\,\mathrm V$、
$\mathcal E_2=1.0\,\mathrm V$
组成的电路，上方电容器
C 两板为 A、B。求稳态电势差
$\varphi_A-\varphi_B$。
\exsolution
稳态电容支路开路，
所以电路主回路只由两电源和三电阻组成。
按图3.47的连接，
用 Kirchhoff 回路电流并分别积分
从 A 到 B 的上下两条路径，
可将开路电容两端电压写为
\[
\boxed{\varphi_A-\varphi_B=
\frac{\mathcal E_2R_3(R_1+R_2)
-\mathcal E_1R_1(R_2+R_3)}
{R_1R_2+R_2R_3+R_3R_1}}.
\]
把给定数值代入，
分母
$10(20)+20(30)+30(10)=1100$，
分子
$1(30)(10+20)-4(10)(20+30)
=900-2000=-1100$，
所以
\[
\boxed{\varphi_A-\varphi_B=-1.0\,\mathrm V}.
\]
\excheck 电容器最终不通直流，
但两端仍可有静态电压；
该电压的符号必须按 A 减 B 计算。

\electrocurrentproblem{185}{118}
图3.48的 Y 形电阻网络，
中心 O 经
$R_1=10\,\Omega$、
$R_2=20\,\Omega$、
$R_3=30\,\Omega$
分别连接外端1、2、3。
三端电势
$\varphi_1=10\,\mathrm V$、
$\varphi_2=6\,\mathrm V$、
$\varphi_3=5\,\mathrm V$。
求流过 $R_1$ 的电流。
\exsolution
把 O 的电势记为 $U$。
由 Kirchhoff 电流定律：
\[
\frac{U-\varphi_1}{R_1}
+\frac{U-\varphi_2}{R_2}
+\frac{U-\varphi_3}{R_3}=0.
\]
所以
\[
U=
\frac{\varphi_1/R_1+\varphi_2/R_2+
\varphi_3/R_3}
{1/R_1+1/R_2+1/R_3}.
\]
从1流向 O 的电流
$I_1=(\varphi_1-U)/R_1$，
消去 $U$ 后得
\[
\boxed{I_1=
\frac{R_3(\varphi_1-\varphi_2)
+R_2(\varphi_1-\varphi_3)}
{R_1R_2+R_2R_3+R_3R_1}}.
\]
代数值：
\[
\boxed{I_1=\frac{30(4)+20(5)}{1100}
=0.20\,\mathrm A}.
\]
\excheck 第1端电势高于中心 O，
故电流从1流向 O，与正结果一致。

\electrocurrentproblem{186}{118}
图3.49：两条横支路自 A 到 B，
上支依次 $R_1=1.0\,\Omega$、
$R_2=2.0\,\Omega$，
下支依次 $R_3=3.0\,\Omega$、
$R_4=4.0\,\Omega$。
两支中点 C、D 以理想导线短接。
$V_A-V_B=25\,\mathrm V$。
求导线 CD 的电流模与方向。
\exsolution
C、D 短接，电势相同为 $U$。
左半的等效电阻
\[
R_L=R_1\parallel R_3
=\frac{1\cdot3}{1+3}=0.75\,\Omega;
\]
右半
\[
R_R=R_2\parallel R_4
=\frac{2\cdot4}{2+4}=\frac43\,\Omega.
\]
总电阻
$R_L+R_R=25/12\,\Omega$，
故总电流
$I_{\rm tot}=25/(25/12)=12\,\mathrm A$。
取 $V_B=0$，
中点电势
$U=I_{\rm tot}R_R=16\,\mathrm V$。
上左流经 $R_1$ 电流
$I_{AC}=(25-16)/1=9$ A，
上右流经 $R_2$ 电流
$I_{CB}=16/2=8$ A。
C 点电荷守恒：
\[
I_{AC}=I_{CB}+I_{CD}
\quad\Longrightarrow\quad
\boxed{I_{CD}=1.0\,\mathrm A}.
\]
方向是
$\boxed{\text{从 C 向 D}}$。
\excheck 不能把短接的 C、D
误当成两个不同电势节点；
差的1 A 通过短线重新分配到下支路。

\electrocurrentproblem{187}{118}
图3.50的五电阻桥式网络，
A—上中节点电阻 $r$、
上中—B 为 $R$、
A—下中为 $R$、
下中—B 为 $r$，
中间桥臂电阻亦为 $r$。
求 A、B 之间的等效电阻。
\exsolution
取 $V_A=V$、$V_B=0$，
上、下中间节点电势分别为 $x,y$。
由节点电流守恒：
\[
\frac{x-V}r+\frac xR+\frac{x-y}r=0,
\]
\[
\frac{y-V}R+\frac yr+\frac{y-x}r=0.
\]
把两式乘 $rR$：
\[
(2R+r)x-Ry=RV,\qquad
-Rx+(2R+r)y=rV.
\]
解线性方程组，代入 A 端注入电流
$I=(V-x)/r+(V-y)/R$，化简为
\[
I=\frac{V(R+3r)}{r(r+3R)}.
\]
所以
\[
\boxed{R_{AB}=\frac VI
=\frac{r(r+3R)}{R+3r}}.
\]
\excheck 特例 $R=r$，
五只相同电阻的对称桥中间无电流，
两条 $2r$ 支路并联为 $r$；
上式给 $r$，相符。

\electrocurrentproblem{188}{118--119}
图3.51：电动势为 $\mathcal E$ 的理想电源
经开关 Sw 与串联底部电阻 $R$ 相连，
在其负载两端还有一只电阻 $R$
与电容器 $C$ 并联。
$t=0$ 闭合开关，电容初始未充电。
求电容电压 $V_C(t)$。
\exsolution
从电容两极端看向外电路，
用戴维南定理化简。
开路电容时，两只 $R$ 串联
由电源 $\mathcal E$ 供电，
端电压
\[
V_{\rm th}=\mathcal E\frac R{R+R}
=\frac{\mathcal E}{2}.
\]
把理想电压源短接时，
从电容端看到两只 $R$ 并联：
\[
R_{\rm th}=R\parallel R=\frac R2.
\]
所以这是初始电压0、时间常数
$\tau=R_{\rm th}C=RC/2$ 的一阶充电过程：
\[
\boxed{V_C(t)=
\frac{\mathcal E}{2}
\left(1-e^{-t/\tau}\right)
=\frac{\mathcal E}{2}
\left(1-e^{-2t/(RC)}\right)}.
\]
\excheck $t=0$ 时电压0，
$t\to\infty$ 时为电阻分压
$\mathcal E/2$；指数衰减时间常数是 $RC/2$，
不是 $RC$。
